\sectioncentered*{Введение}
\addcontentsline{toc}{section}{Введение}

В современном медицинском обслуживании информационные технологии играют ключевую роль в обеспечении качественного и доступного здравоохранения. В этой связи разработка и внедрение интеллектуальных систем, способных обрабатывать и анализировать медицинские данные, становятся всё более актуальными.

Современная медицинская практика сталкивается с рядом вызовов, таких как увеличение объема информации, требующей анализа и обработки, необходимость повышения эффективности диагностики и лечения, а также улучшения качества медицинского обслуживания для пациентов.

Одной из основных проблем, с которой сталкиваются врачи и медицинский персонал, является нехватка времени на анализ медицинских данных, принятие решений и ведение документации. Также часто возникают сложности в точной диагностике и выборе оптимального лечения из-за объема информации и необходимости учета индивидуальных особенностей каждого пациента.

Разработка и внедрение интеллектуальной системы по медицине, такой как медицинский ассистент, могут значительно облегчить задачи медицинского персонала, улучшить точность диагностики, оптимизировать процессы принятия решений и повысить качество медицинского обслуживания для пациентов.

Целью данного проекта является разработка и внедрение интеллектуальной системы по медицине в виде медицинского ассистента, который будет помогать врачам и медицинскому персоналу в анализе и обработке медицинских данных, поддерживая принятие обоснованных и эффективных медицинских решений.

\textbf{Задачи для достижения поставленной цели:}
\begin{itemize}
    \item изучение и анализ предметной области медицины;
    \item создание базы знаний;
    \item проектирование архитектуры и функциональности медицинского ассистента;
    \item реализация программного обеспечения медицинского ассистента;
    \item тестирование и оптимизация работы медицинского ассистента.
\end{itemize}

Важно отметить, что на сегодняшний день многие хотят знать, с
помощью чего эффективнее всего можно контролировать своё здоровье и с помощью каких технологий можно обеспечить достаточно полную диагностику здоровья и реальных актуальных рисков заболеваний, независимо от уровня понимания здоровья человеком.